\section{Algorithms}
\label{sec:algorithm}

\subsection{Problem Setup}
We study cross-node delay prediction from source technology libraries (Nangate45) to target technology libraries (ASAP7).
Given per-arc handcrafted features $\mathbf{x}$ and graph-derived design embedding $\mathbf{z}$, the model predicts normalized delay $y$.
The key challenge is distribution shift across table groups, cell types, and topology patterns.

\subsection{Backbone with Topology-Conditioned Dual Residual (TCDR)}
The shared backbone first produces latent representation
\begin{equation}
\mathbf{h} = f_{\theta}([\mathbf{x};\mathbf{z}]).
\end{equation}
A shared linear regressor gives base prediction:
\begin{equation}
\hat{y}_{\text{base}} = \mathbf{w}^{\top}\mathbf{h} + b.
\end{equation}

To improve OOD generalization, we introduce a topology-conditioned residual expert.
For each sample, topology group ID $t$ is mapped to embedding $\mathbf{e}_t$.
Domain-specific residuals are modeled as
\begin{equation}
\Delta_{\text{src}} = g_{\phi_{\text{src}}}([\mathbf{h};\mathbf{e}_t]),\quad
\Delta_{\text{tgt}} = g_{\phi_{\text{tgt}}}([\mathbf{h};\mathbf{e}_t]).
\end{equation}
Final prediction becomes
\begin{equation}
\hat{y}=
\begin{cases}
\hat{y}_{\text{base}} + \Delta_{\text{src}}, & \text{source domain},\\
\hat{y}_{\text{base}} + \Delta_{\text{tgt}}, & \text{target domain}.
\end{cases}
\end{equation}

Compared with a pure shared head, TCDR decouples universal delay trends (shared base) from topology/domain-specific bias (residual experts).

\subsection{Training Objective}
We jointly optimize source and target losses:
\begin{equation}
\mathcal{L} = \mathcal{L}_{\text{tgt}} + \lambda(e)\,\mathcal{L}_{\text{src}},
\end{equation}
where $\lambda(e)$ is epoch-dependent source weight with annealing.
Both losses use MSE in our final setup.

\subsection{Implementation Notes}
Our strongest and most stable setting in this round uses frozen HGAT encoder, enriched parasitic net features, and TCDR.
Only lightweight residual experts are added on top of the shared head, so the architecture change is minimal and easy to integrate into existing training pipelines.
